\chapter{Weak Convergence and Weak Topology}

\section{Topologies generated by maps}

\begin{notedefinition}
A \emph{topology} $\tau$ on a set $X$ is a family of subsets of $X$ such that
$\varnothing,X\in\tau$, arbitrary unions of members of $\tau$ belong to
$\tau$, and finite intersections of members of $\tau$ belong to $\tau$.  The
pair $(X,\tau)$ is a topological space.

A family $\mathcal B\subseteq\tau$ is a \emph{base} if every open set is a
union of members of $\mathcal B$.  A family $\mathcal S$ is a
\emph{subbase} if the finite intersections of members of $\mathcal S$ form a
base.
\end{notedefinition}

\begin{notedefinition}
A map $F\colon(X,\tau_X)\to(Y,\tau_Y)$ is continuous if
$F^{-1}(O)\in\tau_X$ for every $O\in\tau_Y$.  A sequence $(x_n)$ converges to
$x$ if, for every neighbourhood $U$ of $x$, one has $x_n\in U$ for all
sufficiently large $n$.
\end{notedefinition}

\begin{notedefinition}[Initial topology]
Let $X$ be a set and let $\phi_i\colon X\to(Y_i,\tau_i)$ be maps indexed by
$i\in I$.  The \emph{initial topology} generated by the family $(\phi_i)$ is
the topology with subbase
\[
  \{\phi_i^{-1}(O):i\in I,\ O\in\tau_i\}.
\]
It is the weakest topology on $X$ for which every $\phi_i$ is continuous.
\end{notedefinition}

\begin{noteproposition}
For the initial topology generated by $(\phi_i)$:
\begin{enumerate}
  \item $x_n\to x$ if and only if $\phi_i(x_n)\to\phi_i(x)$ for every $i$;
  \item a map $\psi\colon Z\to X$ is continuous if and only if every
        $\phi_i\circ\psi$ is continuous.
\end{enumerate}
\end{noteproposition}

\begin{proof}
The forward implications follow from continuity of the coordinate maps.  For
the converse in the first statement, a basic neighbourhood of $x$ is a finite
intersection
\[
  \bigcap_{j=1}^m\phi_{i_j}^{-1}(O_j),
  \qquad \phi_{i_j}(x)\in O_j,
\]
and coordinatewise convergence eventually places $x_n$ in every factor.  The
second statement follows because inverse images commute with finite
intersections and arbitrary unions.
\end{proof}

\section{The weak topology}

\begin{notedefinition}
Let $X$ be a normed space.  The \emph{weak topology} $\sigma(X,X^*)$ is the
initial topology generated by all $f\in X^*$.  A sequence converges weakly,
written $x_n\weakto x$, if
\[
  f(x_n)\longrightarrow f(x)
  \qquad(f\in X^*).
\]
\end{notedefinition}

\begin{notetheorem}
The weak topology on a normed space is Hausdorff.
\end{notetheorem}

\begin{proof}
If $x\ne y$, Hahn--Banach gives $f\in X^*$ with $f(x)\ne f(y)$.  Disjoint
open intervals about these two scalar values have disjoint inverse images
containing $x$ and $y$.
\end{proof}

\begin{noteproposition}[Basic weak neighbourhoods]
For $x_0\in X$, the sets
\[
  V(x_0;f_1,\ldots,f_m;\varepsilon)
  =\{x\in X:|f_j(x-x_0)|<\varepsilon,
                    \ 1\le j\le m\}
\]
form a neighbourhood base at $x_0$.  Allowing different positive radii in the
coordinates gives the same topology.
\end{noteproposition}

\begin{proof}
The weak topology is generated by the subbasic sets $f^{-1}(U)$ with
$f\in X^*$ and $U\subseteq\mathbb K$ open.  Hence a basic neighbourhood of
$x_0$ is a finite intersection
\[
  \bigcap_{j=1}^m f_j^{-1}(U_j),
  \qquad f_j(x_0)\in U_j.
\]
For each $j$, choose $\varepsilon_j>0$ such that the open disc of radius
$\varepsilon_j$ about $f_j(x_0)$ is contained in $U_j$.  Then
\[
  \{x:|f_j(x-x_0)|<\varepsilon_j,\quad 1\le j\le m\}
\]
is contained in the given basic neighbourhood.  Conversely, every such
finite-coordinate set is a finite intersection of inverse images of open
discs, so it is weakly open.  Replacing the separate radii by
$\varepsilon=\min_j\varepsilon_j$ proves that one common radius gives the
same neighbourhood base.
\end{proof}

\begin{notetheorem}[Finite-dimensional case]
If $X$ is finite dimensional, weak convergence and norm convergence coincide.
Equivalently, the weak and norm topologies on $X$ are the same.
\end{notetheorem}

\begin{proof}
Choose a basis $e_1,\ldots,e_d$ and its coordinate functionals
$e_1^*,\ldots,e_d^*$.  If $x_n\weakto x$, then every coordinate of $x_n-x$
converges to zero.  All norms on the finite-dimensional coordinate space are
equivalent, so $\|x_n-x\|\to0$.
\end{proof}

\begin{noteexample}
If $(e_n)$ is an orthonormal sequence in a Hilbert space, then
$e_n\weakto0$: Bessel's inequality implies $\inner{x}{e_n}\to0$ for every
$x$.  In contrast, in $\ell^1$ weak convergence of sequences is equivalent to
norm convergence (Schur's theorem).
\end{noteexample}

\begin{noteproposition}[Uniqueness, boundedness, and lower semicontinuity]
Let $X$ be a normed space.
\begin{enumerate}
  \item Weak limits are unique.
  \item If $x_n\weakto x$, then $(x_n)$ is norm bounded.
  \item If $x_n\weakto x$, then
        \[
          \|x\|\le\liminf_{n\to\infty}\|x_n\|.
        \]
\end{enumerate}
\end{noteproposition}

\begin{proof}
Weak limits are unique because the weak topology is Hausdorff.

For boundedness, regard each $x_n$ as the evaluation functional
\[
  J_{x_n}\in(X^*)^*,\qquad J_{x_n}(f)=f(x_n).
\]
For every fixed $f\in X^*$, weak convergence gives
$f(x_n)\to f(x)$, so the scalar sequence $(J_{x_n}(f))$ is bounded.  The
uniform boundedness principle, applied on the Banach space $X^*$, yields
\[
  \sup_n\|J_{x_n}\|<\infty.
\]
Hahn--Banach gives $\|J_{x_n}\|=\|x_n\|$, and hence $(x_n)$ is norm bounded.

For weak lower semicontinuity, if $x=0$ there is nothing to prove.  Otherwise
choose $f\in X^*$ with $\|f\|=1$ and, after multiplying by a unimodular
scalar in the complex case, $f(x)=\|x\|$.  Then
\[
\begin{aligned}
  \|x\|
   &=\operatorname{Re}f(x)
    =\lim_{n\to\infty}\operatorname{Re}f(x_n)\\
   &\le\liminf_{n\to\infty}|f(x_n)|
    \le\liminf_{n\to\infty}\|x_n\|.
\end{aligned}
\]
\end{proof}

\begin{notetheorem}[Weak closure of the unit sphere]
Let $X$ be infinite dimensional and let $S_X=\{x:\|x\|=1\}$.  Then
\[
  \overline{S_X}^{\,w}=B_X=\{x:\|x\|\le1\}.
\]
Consequently, the weak topology is strictly weaker than the norm topology.
\end{notetheorem}

\begin{proof}
Weak lower semicontinuity of the norm shows that
\[
  B_X=\{x:\|x\|\le1\}
\]
is weakly closed.  Since $S_X\subseteq B_X$, this gives
$\overline{S_X}^{\,w}\subseteq B_X$.

For the reverse inclusion, points of $S_X$ already belong to the weak
closure, so let $\|x_0\|<1$.  Consider a basic weak neighbourhood
\[
  V=\{x\in X:|f_j(x-x_0)|<\varepsilon,
                  \ 1\le j\le m\}.
\]
The linear map
\[
  F\colon X\to\mathbb K^m,
  \qquad F(x)=(f_1(x),\ldots,f_m(x)),
\]
has finite-dimensional range.  If its kernel were $\{0\}$, then $X$ would be isomorphic to a subspace of $\mathbb K^m$ and
would be finite dimensional.  Hence
\[
  M:=\bigcap_{j=1}^m\Ker f_j\ne\{0\}.
\]
Choose $0\ne y\in M$.  The function
\[
  h(t)=\|x_0+ty\|,
  \qquad t\ge0,
\]
is continuous, satisfies $h(0)<1$, and tends to infinity as $t\to\infty$.
The intermediate value theorem therefore gives $t_0>0$ with
$\|x_0+t_0y\|=1$.  Since $f_j(y)=0$ for every $j$,
\[
  f_j(x_0+t_0y-x_0)=0,
\]
so $x_0+t_0y\in V\cap S_X$.  Every basic weak neighbourhood of $x_0$ thus
meets the sphere, proving $x_0\in\overline{S_X}^{\,w}$.
\end{proof}

\section{Weak continuity and compactness}

\begin{noteproposition}
Let $X,Y$ be normed spaces.
\begin{enumerate}
  \item If $A\in\Bop{X}{Y}$ and $x_n\weakto x$, then
        $Ax_n\weakto Ax$.
  \item If $A\in\Kop{X}{Y}$ and $x_n\weakto x$, then
        $\|Ax_n-Ax\|\to0$.
  \item If $B\colon X\times Y\to Z$ is a bounded bilinear map, then
        $x_n\to x$ and $y_n\to y$ in norm imply
        $B(x_n,y_n)\to B(x,y)$ in norm.
\end{enumerate}
\end{noteproposition}

\begin{proof}
For the first assertion, $g(Ax_n)=(A^*g)(x_n)\to(A^*g)(x)$ for every
$g\in Y^*$.  For the second, weak convergence makes $(x_n)$ bounded.  If
$Ax_n\not\to Ax$, a subsequence stays a fixed distance away.  Compactness gives
a further norm-convergent subsequence, while the first assertion forces its
weak limit to be $Ax$; hence its norm limit is $Ax$, a contradiction.  The
bilinear estimate
\[
  \|B(x_n,y_n)-B(x,y)\|
  \le \|B\|\bigl(\|x_n-x\|\,\|y_n\|
                  +\|x\|\,\|y_n-y\|\bigr)
\]
proves the third assertion.
\end{proof}

\begin{notetheorem}[Mazur's theorem]
If $x_n\weakto x$ in a normed space, then for every $n$ there is a finite
convex combination of the tail,
\[
  y_n\in\conv\{x_k:k\ge n\},
\]
such that $\|y_n-x\|\to0$.
\end{notetheorem}

\begin{proof}
Let $C_n=\overline{\conv\{x_k:k\ge n\}}^{\|\cdot\|}$.  If $x\notin C_n$,
the geometric Hahn--Banach theorem gives $f\in X^*$ and $\alpha\in\mathbb R$
with
\[
  \operatorname{Re}f(x)>\alpha
  \ge \operatorname{Re}f(z)
  \qquad(z\in C_n).
\]
(Over a real space the real parts may of course be omitted.)  But
$f(x_k)\to f(x)$, contradicting this inequality for all $k\ge n$.  Thus
$x\in C_n$, and one may choose
$y_n\in\conv\{x_k:k\ge n\}$ with $\|y_n-x\|<1/n$.
\end{proof}

\begin{notecorollary}[Banach--Saks--Mazur block form]
If $x_n\weakto x$, one can choose integers
$1\le m_1<m_2<\cdots$ and convex combinations
\[
  z_j\in\conv\{x_k:m_j\le k<m_{j+1}\}
\]
such that $z_j\to x$ in norm.  In particular, weak convergence always admits
strongly convergent convex blocks.
\end{notecorollary}

\begin{proof}
Set $m_1=1$.  Suppose $m_j$ has been chosen.  Mazur's theorem applied to the
tail $(x_k)_{k\ge m_j}$ gives a finite convex combination
\[
  z_j=\sum_{k=m_j}^{N_j}\alpha_k^{(j)}x_k,
  \qquad
  \alpha_k^{(j)}\ge0,
  \quad
  \sum_{k=m_j}^{N_j}\alpha_k^{(j)}=1,
\]
with
\[
  \|z_j-x\|<\frac1j.
\]
Choose $m_{j+1}>N_j$.  Then the blocks are disjoint and
$z_j\in\conv\{x_k:m_j\le k<m_{j+1}\}$.  The displayed estimate gives
$z_j\to x$ in norm.
\end{proof}

\begin{notetheorem}[Convex closures]
For a convex subset $C$ of a normed space,
\[
  \overline C^{\,\|\cdot\|}=\overline C^{\,w}.
\]
Hence a convex set is norm closed if and only if it is weakly closed.
\end{notetheorem}

\begin{proof}
The weak topology is coarser, so norm closure is contained in weak closure.  If
$x\notin\overline C^{\|\cdot\|}$, separate the closed convex set
$\overline C^{\|\cdot\|}$ from $x$ by a continuous linear functional: there
are $f\in X^*$ and $\alpha\in\mathbb R$ such that
\[
  \operatorname{Re}f(x)>\alpha
  \ge \operatorname{Re}f(z)
  \qquad(z\in\overline C^{\|\cdot\|}).
\]
Then
\[
  \{y\in X:\operatorname{Re}f(y)>\alpha\}
\]
is a weakly open neighbourhood of $x$ disjoint from $C$.  Hence $x$ is not in
the weak closure, proving the reverse inclusion.
\end{proof}

\begin{notecorollary}
A convex norm-lower-semicontinuous function
$\varphi\colon X\to(-\infty,\infty]$ is weakly lower semicontinuous.
\end{notecorollary}

\begin{proof}
Every sublevel set $\{x:\varphi(x)\le a\}$ is convex and norm closed, hence
weakly closed.
\end{proof}

\begin{notetheorem}[Weak continuity of linear maps]
Let $E$ and $F$ be Banach spaces and let $T\colon E\to F$ be linear.  Then
$T$ is norm continuous if and only if it is continuous from
$\sigma(E,E^*)$ to $\sigma(F,F^*)$.
\end{notetheorem}

\begin{proof}
Norm continuity implies weak continuity because $f\circ T\in E^*$ for every
$f\in F^*$.  Conversely, weak continuity makes the graph weakly closed in
$E\times F$.  A weakly closed set is norm closed, so the closed graph theorem
implies that $T$ is bounded.
\end{proof}

\section{The Weak-Star Topology}

\begin{notedefinition}
The \emph{weak-$*$ topology} $\sigma(E^*,E)$ on $E^*$ is the initial topology
generated by the evaluation maps
\[
  \phi_x(f)=f(x),\qquad x\in E.
\]
Thus $f_\alpha\weakstarto f$ if and only if
$f_\alpha(x)\to f(x)$ for every $x\in E$.
\end{notedefinition}

\begin{noteproposition}
The weak-$*$ topology is Hausdorff.  A neighbourhood base at $f_0\in E^*$ is
given by
\[
  V(f_0;x_1,\ldots,x_m;\varepsilon)
  =\{f\in E^*:|f(x_j)-f_0(x_j)|<\varepsilon,
                    \ 1\le j\le m\}.
\]
Moreover,
\[
  \text{norm convergence}\Longrightarrow
  \text{weak convergence}\Longrightarrow
  \text{weak-$*$ convergence}
\]
for nets in $E^*$, after identifying the weak topology as
$\sigma(E^*,E^{**})$.
\end{noteproposition}

\begin{proof}
If $f\ne g$, then $f(x)\ne g(x)$ for some $x\in E$.  Disjoint scalar
neighbourhoods of these two values have disjoint inverse images under the
evaluation map $h\mapsto h(x)$, so the weak-$*$ topology is Hausdorff.  Its
neighbourhood base is obtained by taking finite intersections of the
evaluation subbasic sets, exactly as for the weak topology.

If $f_\alpha\to f$ in norm, then for every $F\in E^{**}$,
\[
  |F(f_\alpha-f)|\le\|F\|\,\|f_\alpha-f\|\longrightarrow0,
\]
so norm convergence implies weak convergence.  If $f_\alpha\weakto f$ in
$E^*$, apply the convergence to the canonical functional $J_Ex\in E^{**}$:
\[
  f_\alpha(x)=J_Ex(f_\alpha)
  \longrightarrow J_Ex(f)=f(x)
  \qquad(x\in E).
\]
Thus weak convergence implies weak-$*$ convergence.
\end{proof}

\begin{notelemma}[Finite functional dependence]
Let $\ell,\ell_1,\ldots,\ell_m$ be linear functionals on a vector space $V$.
If
\[
  \bigcap_{j=1}^m\Ker\ell_j\subseteq\Ker\ell,
\]
then $\ell$ belongs to $\Span\{\ell_1,\ldots,\ell_m\}$.
\end{notelemma}

\begin{proof}
Define
\[
  L\colon V\to\mathbb K^m,
  \qquad
  L(v)=(\ell_1(v),\ldots,\ell_m(v)).
\]
If $L(v)=L(w)$, then $v-w\in\bigcap_j\Ker\ell_j\subseteq\Ker\ell$, so
$\ell(v)=\ell(w)$.  Consequently the formula
\[
  \widetilde\ell(L(v))=\ell(v)
\]
defines a well-defined linear functional on the subspace $L(V)$ of
$\mathbb K^m$.  Extend $\widetilde\ell$ to a linear functional on all of
$\mathbb K^m$.  There are scalars $a_1,\ldots,a_m$ such that
\[
  \widetilde\ell(z_1,\ldots,z_m)=\sum_{j=1}^m a_jz_j.
\]
Therefore, for every $v\in V$,
\[
  \ell(v)=\widetilde\ell(L(v))
         =\sum_{j=1}^m a_j\ell_j(v),
\]
which proves $\ell\in\Span\{\ell_1,\ldots,\ell_m\}$.
\end{proof}

\begin{notetheorem}[Dual of the weak-$*$ dual]
A linear functional $\Lambda\colon E^*\to\mathbb K$ is weak-$*$ continuous if
and only if there is an $x\in E$ such that
\[
  \Lambda(f)=f(x)\qquad(f\in E^*).
\]
In other words, the continuous dual of $(E^*,\sigma(E^*,E))$ is the canonical
copy of $E$.
\end{notetheorem}

\begin{proof}
For each $x\in E$, the evaluation map $\phi_x(f)=f(x)$ is weak-$*$
continuous by definition, so every finite linear combination of evaluations
is weak-$*$ continuous.

Conversely, suppose $\Lambda$ is weak-$*$ continuous.  Continuity at the
origin gives $x_1,\ldots,x_m\in E$ and $\varepsilon>0$ such that
\[
  |f(x_j)|<\varepsilon\quad(1\le j\le m)
  \qquad\Longrightarrow\qquad
  |\Lambda(f)|<1.
\]
Let
\[
  M=\bigcap_{j=1}^m\Ker\phi_{x_j}
   =\{f\in E^*:f(x_j)=0\text{ for every }j\}.
\]
If $f\in M$, then $tf$ belongs to the displayed neighbourhood for every
scalar $t$.  Hence
\[
  |t|\,|\Lambda(f)|=|\Lambda(tf)|<1
  \qquad(t\in\mathbb K),
\]
which forces $\Lambda(f)=0$.  Thus $M\subseteq\Ker\Lambda$.  Applying the
finite functional-dependence lemma on the vector space $E^*$ gives scalars
$a_1,\ldots,a_m$ such that
\[
  \Lambda=\sum_{j=1}^m a_j\phi_{x_j}.
\]
Therefore
\[
  \Lambda(f)
  =\sum_{j=1}^m a_j f(x_j)
  =f\left(\sum_{j=1}^m a_jx_j\right),
\]
so $\Lambda$ is evaluation at $x=\sum_j a_jx_j$.
\end{proof}

\begin{notecorollary}
A hyperplane in $E^*$ is weak-$*$ closed if and only if it has the form
\[
  \{f\in E^*:f(x)=a\}
\]
for some $x\in E$ and $a\in\mathbb K$.
\end{notecorollary}

\begin{proof}
Every set of the displayed form is the inverse image of the closed singleton
$\{a\}$ under the weak-$*$ continuous evaluation $f\mapsto f(x)$, and hence
is weak-$*$ closed.

Conversely, write the affine hyperplane as
\[
  H=\{f\in E^*:\Lambda(f)=a\},
\]
where $\Lambda$ is a nonzero algebraic linear functional on $E^*$.  Choose
$f_0\notin H$.  Since $H$ is weak-$*$ closed, there is a balanced weak-$*$
neighbourhood $W$ of $0$ such that
\[
  (f_0+W)\cap H=\varnothing.
\]
Put $d=\Lambda(f_0)-a\ne0$.  We claim that
$|\Lambda(h)|<|d|$ for every $h\in W$.  Otherwise, for some $h\in W$ the
scalar $s=-d/\Lambda(h)$ would satisfy $|s|\le1$.  Balancedness gives
$sh\in W$, while
\[
  \Lambda(f_0+sh)=\Lambda(f_0)-d=a,
\]
contradicting $(f_0+W)\cap H=\varnothing$.  Thus $\Lambda$ is bounded on a
neighbourhood of $0$ and is therefore weak-$*$ continuous.  By the theorem on
the dual of $(E^*,\sigma(E^*,E))$, there is $x\in E$ such that
$\Lambda(f)=f(x)$.  Hence
\[
  H=\{f\in E^*:f(x)=a\},
\]
as required.
\end{proof}

\section{Banach--Alaoglu, Goldstine, and reflexivity}

\begin{notetheorem}[Banach--Alaoglu theorem]
The closed dual unit ball $B_{E^*}$ is compact in the weak-$*$ topology.
\end{notetheorem}

\begin{proof}
For every $x\in E$, let
\[
  K_x=\{z\in\mathbb K:|z|\le\|x\|\},
\]
and form the product
\[
  K=\prod_{x\in E}K_x.
\]
Each $K_x$ is compact, so $K$ is compact by Tychonoff's theorem.  Define
\[
  \Phi\colon B_{E^*}\to K,
  \qquad
  \Phi(f)=(f(x))_{x\in E}.
\]
The map is injective, and the subspace topology induced from the product is
exactly the weak-$*$ topology, because the product coordinate projections
are the evaluation maps $f\mapsto f(x)$.

We show that $\Phi(B_{E^*})$ is closed in $K$.  An element
$\omega=(\omega_x)_{x\in E}\in K$ belongs to the image precisely when it
satisfies
\[
  \omega_{x+y}=\omega_x+\omega_y,
  \qquad
  \omega_{\lambda x}=\lambda\omega_x
\]
for all $x,y\in E$ and scalars $\lambda$.  Each of these is a closed
coordinate condition in the product topology.  Their intersection is
therefore closed.  If $\omega$ satisfies them, the formula
\[
  f_\omega(x)=\omega_x
\]
defines a linear functional, and the fact that $\omega\in K$ gives
\[
  |f_\omega(x)|=|\omega_x|\le\|x\|
  \qquad(x\in E).
\]
Thus $f_\omega\in B_{E^*}$ and $\Phi(f_\omega)=\omega$.  Hence the closed
intersection is exactly $\Phi(B_{E^*})$.

A closed subset of the compact space $K$ is compact, so
$\Phi(B_{E^*})$ is compact.  Since $\Phi$ is a homeomorphism onto its image,
$B_{E^*}$ is weak-$*$ compact.
\end{proof}

\begin{notedefinition}
The canonical embedding $J\colon E\to E^{**}$ is
\[
  (Jx)(f)=f(x).
\]
The space $E$ is \emph{reflexive} if $J(E)=E^{**}$.
\end{notedefinition}

\begin{notetheorem}[Goldstine theorem]
The canonical image $J(B_E)$ is weak-$*$ dense in $B_{E^{**}}$.
\end{notetheorem}

\begin{proof}
Fix $x^{**}\in B_{E^{**}}$.  A basic weak-$*$ neighbourhood of $x^{**}$ is
determined by $f_1,\ldots,f_m\in E^*$ and $\varepsilon>0$; we must find
$x\in B_E$ such that
\[
  |f_j(x)-x^{**}(f_j)|<\varepsilon
  \qquad(1\le j\le m).
\]
Define
\[
  Q\colon E\to\mathbb K^m,
  \qquad
  Qx=(f_1(x),\ldots,f_m(x)),
\]
and put
\[
  y=(x^{**}(f_1),\ldots,x^{**}(f_m)).
\]
We claim that $y\in\overline{Q(B_E)}$.

Suppose not.  The closed convex set $\overline{Q(B_E)}$ and the point $y$
can be strictly separated in the finite-dimensional real vector space
underlying $\mathbb K^m$.  Thus there is a real-linear functional of the form
\[
  z\longmapsto\operatorname{Re}\sum_{j=1}^m a_jz_j
\]
for which
\[
  \operatorname{Re}\sum_{j=1}^m a_jx^{**}(f_j)
  >\sup_{x\in B_E}\operatorname{Re}\sum_{j=1}^m a_jf_j(x).
\]
Let $f=\sum_j a_jf_j$.  The left side is
$\operatorname{Re}x^{**}(f)$, while the right side is $\|f\|$: after
multiplying a nearly norming vector by a unimodular scalar, the real part can
be made arbitrarily close to $\|f\|$.  We would therefore have
\[
  \operatorname{Re}x^{**}(f)>\|f\|,
\]
contradicting
$|x^{**}(f)|\le\|x^{**}\|\,\|f\|\le\|f\|$.  Hence
$y\in\overline{Q(B_E)}$.

Choose $x\in B_E$ with $\|Qx-y\|_\infty<\varepsilon$.  Then the displayed
coordinate inequalities hold, so the given weak-$*$ neighbourhood of
$x^{**}$ meets $J(B_E)$.  This proves weak-$*$ density.
\end{proof}

\begin{notetheorem}[Kakutani's characterization]
A Banach space $E$ is reflexive if and only if its closed unit ball is weakly
compact.
\end{notetheorem}

\begin{proof}
If $E$ is reflexive, $J$ identifies $B_E$ with $B_{E^{**}}$, which is weak-$*$
compact by Banach--Alaoglu; under $J$ this is the weak topology of $E$.
Conversely, if $B_E$ is weakly compact, then $J(B_E)$ is weak-$*$ compact and
hence weak-$*$ closed in $E^{**}$.  Goldstine says it is weak-$*$ dense in
$B_{E^{**}}$, so $J(B_E)=B_{E^{**}}$.  Scaling gives $J(E)=E^{**}$.
\end{proof}

\begin{notetheorem}[Eberlein--\v{S}mulian theorem]
A subset of a Banach space is relatively weakly compact if and only if every
sequence in it has a weakly convergent subsequence whose limit belongs to its
weak closure.  In particular, a Banach space $E$ is reflexive if and only if
$B_E$ is weakly sequentially compact.
\end{notetheorem}

\begin{proof}[The direction proved by the separable-reduction idea in the source]
Assume that $E$ is reflexive and let $(x_n)$ be a sequence in $B_E$.  Set
\[
  M_0=\Span\{x_n:n\ge1\},
  \qquad
  M=\overline{M_0}^{\|\cdot\|}.
\]
Then $M$ is separable.  It is also reflexive, because a closed subspace of a
reflexive Banach space is reflexive.  Hence $B_M$ is weakly compact by
Kakutani's theorem.

To recover the metrizability step used in the handwritten proof, note first
that separability of $M$ makes $B_{M^*}$ weak-$*$ compact and metrizable.
Because $M$ is reflexive, the weak and weak-$*$ topologies on $B_{M^*}$
coincide.  Thus $B_{M^*}$ is a compact metric space and has a countable
weakly dense subset $(g_j)$.  The norm-closed convex hull of $(g_j)$ is all
of $B_{M^*}$: otherwise Hahn--Banach would separate a point of $B_{M^*}$
from that convex hull by an element of $M^{**}=M$, which is a weakly
continuous functional and would contradict weak density.  Therefore $M^*$
is norm separable.

Choose a norm-dense sequence $(f_j)$ in $B_{M^*}$.  On $B_M$, the metric
\[
  d(u,v)=\sum_{j=1}^{\infty}2^{-j}
         \frac{|f_j(u-v)|}{1+|f_j(u-v)|}
\]
induces the weak topology.  Hence the weakly compact set $B_M$ is compact
and metrizable, and therefore sequentially compact.  The sequence $(x_n)$
thus has a subsequence converging weakly in $M$, and hence weakly in $E$.
This proves the implication
\[
  E\text{ reflexive}\quad\Longrightarrow\quad
  B_E\text{ weakly sequentially compact}
\]
by exactly the separable-subspace strategy of the notes.

The converse implication and the general subset statement constitute the
deep half of the Eberlein--\v{S}mulian theorem; the source does not contain
the additional diagonal and closure arguments needed for a complete proof of
that half.
\end{proof}

\begin{notecorollary}
If $E$ is reflexive, every closed bounded convex subset of $E$ is weakly
compact.
\end{notecorollary}

\begin{proof}
Such a set lies in a scalar multiple of $B_E$, which is weakly compact.  It is
weakly closed because convex norm-closed sets are weakly closed.
\end{proof}

\section{Separability and metrizability}

\begin{notetheorem}
If $E^*$ is separable, then $E$ is separable.
\end{notetheorem}

\begin{proof}
Choose a norm-dense sequence $(f_n)$ in $S_{E^*}$.  For each $n$, choose
$x_n\in S_E$ with $|f_n(x_n)|>1/2$.  Let
$M=\overline{\Span\{x_n:n\ge1\}}$.  If $M\ne E$, Hahn--Banach gives
$f\in S_{E^*}$ with $f|_M=0$.  Choose $f_n$ with $\|f-f_n\|<1/2$.  Then
\[
  |f_n(x_n)|=|(f_n-f)(x_n)|<1/2,
\]
a contradiction.  Thus $M=E$, and rational linear combinations of the
$x_n$ form a countable dense set.
\end{proof}

\begin{notetheorem}
For a Banach space $E$, the following are equivalent:
\begin{enumerate}
  \item $E$ is separable;
  \item $B_{E^*}$ is metrizable in the weak-$*$ topology.
\end{enumerate}
\end{notetheorem}

\begin{proof}
Assume first that $E$ is separable and choose a norm-dense sequence
$(x_n)$ in $B_E$.  Define, for $f,g\in B_{E^*}$,
\[
  d(f,g)=\sum_{n=1}^{\infty}2^{-n}
         \frac{|(f-g)(x_n)|}{1+|(f-g)(x_n)|}.
\]
This is a metric.  Indeed, if $d(f,g)=0$, then $f$ and $g$ agree on every
$x_n$, hence by continuity on $B_E$ and therefore on all of $E$.

Weak-$*$ convergence implies convergence for $d$, because each summand tends
to zero and the summands are dominated by $2^{-n}$; a finite-head plus
small-tail argument proves convergence of the series.  Conversely, suppose
$d(f_k,f)\to0$ with $f_k,f\in B_{E^*}$.  Then
$f_k(x_n)\to f(x_n)$ for every $n$.  Fix $x\in E$ and $\varepsilon>0$.
For $x\ne0$, choose $x_n$ such that
\[
  \left\|\frac{x}{\|x\|}-x_n\right\|
  <\frac{\varepsilon}{4\|x\|};
\]
for $x=0$ there is nothing to prove.  Since
$\|f_k-f\|\le2$,
\[
\begin{aligned}
  |(f_k-f)(x)|
   &\le |(f_k-f)(x-\|x\|x_n)|
      +\|x\|\,|(f_k-f)(x_n)|\\
   &\le 2\|x-\|x\|x_n\|
      +\|x\|\,|(f_k-f)(x_n)|.
\end{aligned}
\]
The first term is $<\varepsilon/2$, and the second is eventually
$<\varepsilon/2$.  Hence $f_k(x)\to f(x)$ for every $x\in E$.  Thus $d$
induces the weak-$*$ topology on $B_{E^*}$.

Conversely, suppose $B_{E^*}$ is weak-$*$ metrizable.  Banach--Alaoglu makes
it a compact metric space.  Therefore $C(B_{E^*})$, with the supremum norm,
is separable.  Define
\[
  J\colon E\to C(B_{E^*}),
  \qquad
  (Jx)(f)=f(x).
\]
Weak-$*$ continuity of each evaluation shows $Jx\in C(B_{E^*})$, and
Hahn--Banach gives
\[
  \|Jx\|_\infty
  =\sup_{\|f\|\le1}|f(x)|
  =\|x\|.
\]
Thus $J$ is an isometric embedding.  Every subspace of a separable metric
space is separable, so $E$ is separable.
\end{proof}

\begin{notetheorem}
The weak topology on $B_E$ is metrizable if and only if $E^*$ is separable.
\end{notetheorem}

\begin{proof}
Suppose first that $E^*$ is separable and choose a norm-dense sequence
$(f_n)$ in $B_{E^*}$.  Define
\[
  d(x,y)=\sum_{n=1}^{\infty}2^{-n}
         \frac{|f_n(x-y)|}{1+|f_n(x-y)|},
  \qquad x,y\in B_E.
\]
The function $d$ is a metric because the dense family $(f_n)$ separates
points.  Weak convergence in $B_E$ implies $d$-convergence by dominated
convergence for the series.  Conversely, if $d(x_k,x)\to0$, then
$f_n(x_k-x)\to0$ for every $n$.  For arbitrary $f\in B_{E^*}$ and
$\varepsilon>0$, choose $f_n$ with $\|f-f_n\|<\varepsilon/4$.  Since
$x_k,x\in B_E$,
\[
\begin{aligned}
  |f(x_k-x)|
   &\le |(f-f_n)(x_k-x)|+|f_n(x_k-x)|\\
   &\le2\|f-f_n\|+|f_n(x_k-x)|,
\end{aligned}
\]
which is eventually $<\varepsilon$.  By homogeneity the same holds for every
$f\in E^*$, so $x_k\weakto x$.  Thus $d$ induces the weak topology on
$B_E$.

Conversely, suppose the relative weak topology on $B_E$ is metrizable.  It
is then first countable at $0$, so choose a countable neighbourhood base
$(U_n)$ at $0$ in $B_E$.  For each $n$, choose a basic weak neighbourhood
\[
  V_n=\{x\in E:|f(x)|<\varepsilon_n
                 \text{ for every }f\in F_n\}
\]
with $F_n\subseteq E^*$ finite and
\[
  0\in B_E\cap V_n\subseteq U_n.
\]
Let
\[
  \mathcal F=\bigcup_{n=1}^{\infty}F_n,
  \qquad
  M=\overline{\Span\mathcal F}^{\|\cdot\|}\subseteq E^*.
\]
We prove $M=E^*$.  Otherwise choose $g\in E^*\setminus M$.  Hahn--Banach,
applied in the Banach space $E^*$, gives $x^{**}\in E^{**}$ such that
\[
  x^{**}|_M=0,
  \qquad
  x^{**}(g)\ne0.
\]
After normalization, assume $\|x^{**}\|\le1$.  Goldstine's theorem gives a
net $(x_\alpha)\subseteq B_E$ such that
\[
  Jx_\alpha\weakstarto x^{**}
  \quad\text{in }E^{**}.
\]
For every $f\in\mathcal F$, we have
$f(x_\alpha)\to x^{**}(f)=0$.  Hence, for each $n$, the net is eventually in
$B_E\cap V_n\subseteq U_n$.  Since $(U_n)$ is a neighbourhood base, this
means $x_\alpha\weakto0$ in $E$.  In particular, $g(x_\alpha)\to0$.
Goldstine convergence, however, also gives
\[
  g(x_\alpha)=(Jx_\alpha)(g)\longrightarrow x^{**}(g)\ne0,
\]
a contradiction.  Thus $M=E^*$, and the countable set $\mathcal F$ has norm-
dense linear span.  Therefore $E^*$ is separable.
\end{proof}

\begin{notecorollary}
If $E$ is separable, then every bounded sequence in $E^*$ has a weak-$*$
convergent subsequence.
\end{notecorollary}

\begin{proof}
After scaling, the sequence lies in $B_{E^*}$, which is weak-$*$ compact and
metrizable.  Compact metric spaces are sequentially compact.
\end{proof}

\section{Uniform convexity and reflexivity}

\begin{notedefinition}
A Banach space $E$ is \emph{uniformly convex} if, for every $\varepsilon>0$,
there is $\delta>0$ such that
\[
  \|x\|,\|y\|\le1,
  \quad \|x-y\|\ge\varepsilon
  \quad\Longrightarrow\quad
  \left\|\frac{x+y}{2}\right\|\le1-\delta.
\]
\end{notedefinition}

\begin{noteproposition}[Radon--Riesz property]
Let $E$ be uniformly convex.  If $x_n\weakto x$ and
$\|x_n\|\to\|x\|$, then $x_n\to x$ in norm.
\end{noteproposition}

\begin{proof}
The case $x=0$ is immediate.  Otherwise normalize
$y_n=x_n/\|x_n\|$ and $y=x/\|x\|$.  Then $y_n\weakto y$ and
$\|y_n\|=\|y\|=1$.  If a subsequence satisfied
$\|y_n-y\|\ge\varepsilon$, uniform convexity would give
$\|(y_n+y)/2\|\le1-\delta$.  But weak lower semicontinuity applied to
$(y_n+y)/2\weakto y$ gives
$1\le\liminf\|(y_n+y)/2\|$, a contradiction.
\end{proof}

\begin{notetheorem}[Milman--Pettis theorem]
Every uniformly convex Banach space is reflexive.
\end{notetheorem}

\begin{proof}
Let $x^{**}\in S_{E^{**}}$.  By Goldstine there is a net
$(x_\alpha)\subseteq B_E$ with $Jx_\alpha\weakstarto x^{**}$.  Fix
$\varepsilon>0$ and choose the corresponding uniform-convexity constant
$\delta>0$.  After multiplying a norming functional by a unimodular scalar, choose
$f\in S_{E^*}$ with
$\operatorname{Re}x^{**}(f)>1-\delta/2$.  Eventually
$\operatorname{Re}f(x_\alpha)>1-\delta$, and hence for
large $\alpha,\beta$,
\[
  \left\|\frac{x_\alpha+x_\beta}{2}\right\|
  \ge \operatorname{Re}f\!\left(\frac{x_\alpha+x_\beta}{2}\right)
  >1-\delta.
\]
Uniform convexity gives $\|x_\alpha-x_\beta\|<\varepsilon$.  Thus the net is
norm Cauchy and converges to some $x\in B_E$.  Since norm convergence implies
weak-$*$ convergence under $J$, uniqueness of the weak-$*$ limit gives
$Jx=x^{**}$.  Scaling proves that $J$ is surjective.
\end{proof}
